\documentclass[10pt, a4paper]{article}

% --- Packages ---
\usepackage[a4paper, left=2cm, right=2cm, top=2cm, bottom=2cm]{geometry} % Reduced margins to 2cm for more text space
\usepackage{times} % Times New Roman font
\usepackage{fancyhdr}
\usepackage{amsmath,amssymb,amsthm}
\usepackage{graphicx}
\usepackage{array}
\usepackage{setspace}
\usepackage{hyperref}

% --- Custom Command Definitions from thesis ---
\newcommand{\f}[1]{\textit{#1}}
\newcommand{\bi}[1]{\textbf{\textit{#1}}}
\newcommand{\bo}[1]{\textbf{#1}}

% --- Page Style & Header/Footer Configuration ---
\pagestyle{fancy}
\fancyhf{}
\fancyhead[L]{\small Hanafi \& Widyaningsih}
\fancyhead[R]{\small Pendekatan STATIS dan Dynamic Biplot untuk Analisis Spasio-Temporal\dots}
\fancyfoot[C]{\small Prosiding Diskusi Panel Nasional Pendidikan Matematika Universitas Indraprasta PGRI Jakarta.}
\renewcommand{\headrulewidth}{0.4pt}
\renewcommand{\footrulewidth}{0.4pt}

% --- Paragraph Indentation & Spacing ---
\setlength{\parindent}{0cm}
\singlespacing

\begin{document}

% --- Category Header ---
{\itshape Original Research / Literature Review / Book Review / Lecture Note}
\vspace{0.8em}

% --- Title Block ---
\begin{flushright}
    {\fontsize{15}{18}\selectfont\bfseries Pendekatan STATIS dan Dynamic Biplot untuk Analisis\\ Spasio-Temporal Sosio-Ekonomi\\ (Studi Kasus: Provinsi Banten, DKI Jakarta, dan Jawa Barat 2020--2024)\par}
    \vspace{0.8em}
    {\large\bfseries Ammar Hanafi\textsuperscript{1,*)} \& Yekti Widyaningsih\textsuperscript{1)}}\par
    \vspace{0.3em}
    {\normalsize\itshape \textsuperscript{1}Program Studi Sarjana Statistika, Departemen Matematika,\\ Fakultas Matematika dan Ilmu Pengetahuan Alam,\\ Universitas Indonesia, Depok 16424, Indonesia\par}
\end{flushright}

\vspace{0.8em}

% --- Info Articles Table (Single-page standard tabular, compact 9pt font) ---
\noindent
\begingroup
\small % Set font to 9pt inside table for tight fit
\begin{tabular}{@{} p{0.28\textwidth} @{} p{0.04\textwidth} @{} p{0.68\textwidth} @{}}
\hline
\rule{0pt}{2.5ex}% spacing
\textbf{INFO ARTICLES} & & \\
\cline{1-1}
\textbf{Key Words:} \par \vspace{0.4em}
Analisis Tiga-arah; Data Longitudinal; K-Means Clustering; PCA; Visualisasi Interaktif \par \vspace{2em}
This article is licensed under a Creative Commons Attribution-ShareAlike 4.0 International License.
& &
\textbf{Abstract} \par \vspace{0.4em}
This study analyzes the spatio-temporal socio-economic dynamics across 41 regencies and cities in Banten, DKI Jakarta, and West Java from 2020 to 2024. Addressing the limitations of traditional multivariate methods during pandemic contractions and recoveries, the STATIS method is applied to analyze 12 strategic socio-economic indicators. The method constructs a stable compromise structure ($\mathbf{W}_c$) from Spectral-Decomposition-weighted annual matrices to map dominant regional development patterns. Spatio-temporal trajectories and klaster shifts are visualized in a low-dimensional space using the Dynamic Biplot and K-Means clustering. The entire pipeline is implemented in an interactive Python web dashboard, integrating analytical modules from diagnostic tests to transition matrices. The results reveal strong longitudinal structural consistency (RV > 0.90), distinguishing metropolitan V-shaped recoveries from southern stagnant trajectories. This system provides policymakers with an advanced analytical tool to monitor regional resilience and design targeted development strategies. \par \vspace{1.5em}
\textbf{Abstrak} \par \vspace{0.4em}
Penelitian ini menganalisis dinamika sosio-ekonomi spasio-temporal di 41 kabupaten/kota Provinsi Banten, DKI Jakarta, dan Jawa Barat periode 2020--2024. Guna mengatasi keterbatasan metode multivariat tradisional, metode STATIS diterapkan untuk mengintegrasikan 12 indikator pembangunan strategis. Melalui dekomposisi spektral matriks tahunan, dibentuk struktur kompromi ($\mathbf{W}_c$) yang menggambarkan pola pembangunan bersama yang stabil. Visualisasi trajektori spasio-temporal dan pergeseran klaster K-Means dalam ruang berdimensi rendah dilakukan dengan pendekatan \textit{Dynamic Biplot}. Seluruh prosedur analisis diimplementasikan ke dalam aplikasi \textit{dashboard} web interaktif berbasis Python yang mengintegrasikan pengujian diagnostik hingga matriks transisi. Hasil analisis menunjukkan konsistensi struktural longitudinal yang tinggi ($RV > 0,90$), membedakan pemulihan berbentuk-V wilayah metropolitan dari trajektori stagnan di wilayah selatan. Sistem ini menjadi instrumen analitik canggih bagi pembuat kebijakan untuk memantau resiliensi dan merancang strategi pemerataan pembangunan regional. \\
\cline{1-1}
\multicolumn{3}{@{} p{\textwidth} @{}}{%
\rule{0pt}{2.5ex}\textbf{Correspondence Address:} Program Studi Sarjana Statistika, Departemen Matematika, Fakultas Matematika dan Ilmu Pengetahuan Alam, Universitas Indonesia, Depok 16424, Indonesia. e-mail: \href{mailto:ammar.hanafi@ui.ac.id}{ammar.hanafi@ui.ac.id} \newline
\textbf{How to Cite (APA 6th Style):} Hanafi, A. (2026). Pendekatan STATIS dan Dynamic Biplot untuk Analisis Spasio-Temporal Sosio-Ekonomi (Studi Kasus: Provinsi Banten, DKI Jakarta, dan Jawa Barat 2020--2024). \textit{Prosiding Diskusi Panel Nasional Pendidikan Matematika}, page-page. \newline
\textbf{Copyright:} Author(s), (2026)
} \\
\hline
\end{tabular}
\endgroup

\end{document}
